% Teil II: der durchgerechnete Weg. Reihenfolge beibehalten.

\subsection*{Gegeben und gesucht}
Vier Grenzwerte von Brüchen bzw.\ unbestimmten Formen, z. B.\ bei $0/0$ oder
$\infty/\infty$. Gesucht: die jeweiligen Grenzwerte.

\subsection*{Ansatz}
In jedem Teil zuerst prüfen, welche unbestimmte Form vorliegt. Dann Zähler
und Nenner getrennt ableiten und den neuen Bruch auswerten; bei Bedarf die
Regel wiederholen. Typische Resultate dabei sind bekannte Grenzwerte wie
$\lim_{u\to0}\frac{\sin u}{u}=1$ oder $\lim_{u\to0}\frac{\mathrm e^u-1}{u}=1$.

\subsection*{Rechnung (Methode)}
Für einen Bruch $\lim\frac{f}{g}$ mit $0/0$ oder $\infty/\infty$:
\begin{equation}
  \lim_{x\to x_0}\frac{f(x)}{g(x)}
  =\lim_{x\to x_0}\frac{f'(x)}{g'(x)},
  \qquad
  \text{ggf.\ erneut anwenden.}
\end{equation}
Für Formen wie $0\cdot\infty$ oder $1^\infty$ zuerst umformen
(Kehrwert, Logarithmus), dann ableiten.

\begin{achtung}[Genaue Angaben im Original prüfen]
Die vier konkreten Grenzwerte auf dem Aufgabenblatt sind im eingescannten
Bild für das Textlesen schlecht erkennbar. Die Methode oben ist die prüfbare
Vorgehensweise; die konkreten Brüche bitte am Blatt ablesen und damit
nachrechnen, statt die (nicht sicher lesbaren) Ausdrücke aus diesem Hilfsblatt
zu übernehmen.
\end{achtung}

\begin{merksatz}
Vor jedem l'Hôpital-Schritt die unbestimmte Form prüfen, dann Zähler und
Nenner getrennt ableiten. Andere unbestimmte Formen erst auf $0/0$ oder
$\infty/\infty$ bringen.
\end{merksatz}
